../cictikz/src/cictikz/data/tex/cictikz_preamble.tex

% The current mode DAC of dac_i again, with the switch drive tied down.
%
% V_bias is a constant from outside. It gates the cascode in the reference
% branch and runs the length of the array as a rail, and one side of every
% steering pair sits on it while the other side takes the bit. The bit
% therefore has to cross V_bias in both directions for the pair to tip, and
% only has to move far enough either side of it to do so: a small swing about
% a fixed level disturbs the tail node — and so the settling — far less than
% a full swing does. The inset says that in one picture.
%
% In the reference branch the diode connection is taken from the top of the
% stack, above the cascode, so the mirror's gate carries the voltage of the
% whole stack rather than of its own drain. That node is the gate rail the
% cells' tail devices share. Both reference devices are mirrored so their
% gates face the array.

\begin{circuitikz}[american, thick, transform shape, circuit ee IEC]

  ../cictikz/src/cictikz/data/tex/cictikz_lib.tex
  ../cictikz/src/cictikz/data/tex/rdac_lib.tex

  \def\ytail{0.6}       % source of the tail devices
  \def\ygate{1.4}       % gate rail of the tail devices
  \def\ybias{2.7}       % the V_bias rail
  \def\ypair{3.4}       % source of the steering pairs
  \def\ydrain{5.2}      % drain of the steering pairs
  \def\ydrv{4.3}        % the drivers, on the gates of the pair
  \def\yplus{6.1}       % dummy rail, into the + input
  \def\yminus{6.9}      % summing rail, into the - input
  \def\xref{0}
  \def\dw{0.85}         % half the spacing of a steering pair
  \def\cw{3.0}          % half the width of a cell
  \def\xa{5.6}          % the N cell
  \def\xg{9.3}          % the cells that are not drawn
  \def\xb{13.0}         % the 1 cell

  % ---- Reference branch: a current source into a cascoded mirror ----
  \draw (\xref,\ytail) to [Tnmos, n=M0, mirror] (\xref,{\ytail+1.6});
  \draw (\xref,\ytail) -- ++(0,-0.3) \vground;
  \draw (\xref,{\ytail+1.6}) to [Tnmos, n=MC, mirror] (\xref,{\ytail+3.2});
  \draw (\xref,{\ytail+3.2}) to [I, invert] (\xref,{\ytail+5.0});
  \draw (\xref,{\ytail+5.0}) \vsupply;
  % The diode connection is taken from the top of the stack, above the
  % cascode, so the mirror's gate sits at the whole stack's voltage rather
  % than at its own drain. That node is the gate rail the cells share.
  \draw (M0.gate) -- (1.7,\ygate) -- (1.7,{\ytail+3.2}) -- (\xref,{\ytail+3.2});
  \fill (\xref,{\ytail+3.2}) circle (0.07);
  \fill (1.7,\ygate) circle (0.07);
  % V_bias comes from outside and gates the cascode. It is the same constant
  % the fixed side of every steering pair sits on, so it leaves as a rail.
  \draw (MC.gate) -- (1.0,{\ytail+2.4}) -- (1.0,\ybias);
  \draw (1.35,\ybias) -- (1.35,2.15) coordinate (vbport);
  \fill (1.35,\ybias) circle (0.07);
  \rdacPort{(vbport)}{north east}{$V_{bias}$}

  % ---- One DAC cell, drawn once ----
  % #1 x origin, #2 suffix for the node names, #3 the device size label
  \newcommand{\dacCell}[3]{%
    \draw (#1,\ytail) to [Tnmos, n=MT#2] (#1,{\ytail+1.6});
    \draw (#1,\ytail) -- ++(0,-0.3) \vground;
    \draw (#1,{\ytail+1.6}) -- (#1,\ypair);
    % The size marking sits beside its own device, clear of both rails.
    \node[anchor=west] at ({#1+0.3},{\ytail+1.25}) {#3};
    \draw ({#1-\dw},\ypair) -- ({#1+\dw},\ypair);
    \draw ({#1-\dw},\ypair) to [Tnmos, n=ML#2] ({#1-\dw},\ydrain);
    \draw ({#1+\dw},\ypair) to [Tnmos, n=MR#2, mirror] ({#1+\dw},\ydrain);
    %
    % One side of the pair sits on the constant V_bias rail, the other takes
    % the bit. The bit swings about V_bias, so the pair tips one way or the
    % other and hands the cell current between the two rails.
    \draw (ML#2.gate) -- ({#1-\dw-0.55},\ydrv) -- ({#1-\dw-0.55},\ybias);
    \fill ({#1-\dw-0.55},\ybias) circle (0.07);
    \draw (MR#2.gate) -- ({#1+\cw},\ydrv);
    \node[anchor=south] at ({#1+\cw-0.55},{\ydrv+0.16}) {$\overline{b_{#2}}$};
    %
    % b steers into the summing rail, b-bar into the dummy rail.
    \draw ({#1-\dw},\ydrain) -- ({#1-\dw},\yminus);
    \fill ({#1-\dw},\yminus) circle (0.07);
    \draw ({#1+\dw},\ydrain) -- ({#1+\dw},\yplus);
    \fill ({#1+\dw},\yplus) circle (0.07);
    \draw[dashed, thin] ({#1-\cw},{\ytail-0.75}) rectangle ({#1+\cw},{\ydrain+0.55});
  }

  % The cells that are not drawn: the same boundary, left blank and narrow.
  \newcommand{\dacGhost}[1]{%
    \draw[dashed, thin] ({#1-0.7},{\ytail-0.75}) rectangle ({#1+0.7},{\ydrain+0.55});
    \node at (#1,{0.5*(\ytail+\ydrain)}) {$\cdots$};
  }

  \dacCell{\xa}{n}{$N$}
  \dacGhost{\xg}
  \dacCell{\xb}{0}{$1$}

  \node[anchor=north, align=center] at (\xg,{\ytail-1.05})
    {$N$ copies of the same cell, sized $1, 2, 4 \ldots N$};

  % ---- The two rails out of the reference branch ----
  \draw (1.7,\ygate) -- (MTn.gate);
  \draw (MTn.gate) -- (MT0.gate);
  \fill (MTn.gate) circle (0.07);
  \draw (1.0,\ybias) -- ({\xb+\cw},\ybias);

  % ---- The two rails into the amplifier ----
  \draw ({\xa-\dw},\yminus) -- (17.8,\yminus);
  \draw ({\xa+\dw},\yplus) -- (17.8,\yplus);
  \rdacAmp{OA}{(17.8,6.5)}
  \draw (17.8,\yminus) -- (OAinn);
  \draw (17.8,\yplus) -- (OAinp);

  % ---- The dummy rail is held at a reference voltage ----
  \draw (16.8,\yplus) to [V] (16.8,{\yplus-2.2});
  \draw (16.8,{\yplus-2.2}) -- ++(0,-0.3) \vground;
  \fill (16.8,\yplus) circle (0.07);

  % ---- Feedback resistor over the top ----
  \draw (17.3,\yminus) -- (17.3,9.3) -- ++(0.7,0) \hresistor{$R_F$} -- ++(0.7,0)
        -- (21.0,9.3) -- (21.0,6.5) -- (OAout);
  \fill (17.3,\yminus) circle (0.07);
  \draw (OAout) -- ++(0.9,0);

  % ---- Inset: what V_bias is for ----
  \begin{scope}[shift={(6.0,8.6)}]
    \draw[dashed] (0,0.55) -- (4.4,0.55) node[anchor=west] {$V_{bias}$};
    \draw[thick] (0,1.15) -- (1.1,1.15) -- (1.1,-0.05) -- (2.2,-0.05)
                 -- (2.2,1.15) -- (3.3,1.15) -- (3.3,-0.05) -- (4.4,-0.05);
    \node[anchor=east] at (-0.15,0.55) {$\overline{b_n}$};
    \node[anchor=north, align=center] at (2.2,-0.35)
      {the drive swings about $V_{bias}$, not rail to rail};
  \end{scope}

\end{circuitikz}
\end{document}
